\documentclass{article}
\usepackage[utf8]{inputenc}
\usepackage{geometry}
\usepackage{enumitem}
\usepackage{hyperref}
\usepackage{xcolor}
\usepackage{graphicx}
\usepackage{array}
\usepackage{multirow}
\usepackage{amsmath, amssymb}
\usepackage{chemfig}
\usepackage{mhchem}
\usepackage{tikz}
\usepackage{setspace}
\usepackage{soul}
\usepackage{mathtools}
\usepackage{booktabs}
\usepackage{chemformula}
\usepackage{siunitx}
\usepackage{longtable}
\geometry{margin=1in}
\setlength{\parskip}{0.5em}

\begin{document}

\usetikzlibrary{positioning, shapes.geometric, arrows.meta, decorations.pathreplacing, fit, backgrounds}
\begin{center}
    {\LARGE \textbf{SI Session Planning Form}}
\end{center}

\vspace{0.5cm}

\noindent
\begin{flushleft}
\textbf{SI Leader:} Brandon Cobb \\
\textbf{Session Week:} 2 \\
\textbf{Session \#:} 2 \\
\textbf{Instructor:} Professor Dudgeon \\
\textbf{Old Plans:} \href{https://macombedu.sharepoint.com/:b:/s/ASC-SupplementalInstructionEmbeddedTutoring/IQBvRyfNvxGGQYObfgZOKn8UARbteAbUZFaLhS7nR4OoDmQ?e=Ic6HHf}{SharePoint} \\
\textbf{Session Goal:} Students will be able to take measurements, perform dimensional analysis, appropriately round to the correct significant figures, count significant figures, determine exact and inexact numbers and understand scientific notation.
\end{flushleft}

\vspace{1em}
\section*{\underline{Opener}}
{\centering\large\textbf{\hl{Take Attendance}}\par}\vspace{-0.25\baselineskip}


\begin{tabular}{|p{4.2cm}|p{9.8cm}|}
\hline
\textbf{Collaborative Learning Techniques} & Individual Presentation \\
\hline
\textbf{Learning Strategy} & I'm a mathematician! \\
\hline
\textbf{Time} & 12:00---12:10 \\
\hline
\textbf{Materials} & Ready to speak \\
\hline
\textbf{Lecture/Textbook/Old Plans/Slide References} & Computer, class roster, Chapter 2: Measurments \\
\hline
\end{tabular}

\vspace{0.5cm}
\noindent
{\large \textbf{Instructions (please provide a \hl{DETAILED} activity breakdown below (and/or attached}} \\
\newpage


\vspace{0.5cm}
\noindent
{\large \textbf{Instructions (please provide a \hl{DETAILED} activity breakdown below (and/or attached}} \\

\begin{tabular}{|p{0.9\linewidth}|}
\hline

\vspace{0.2cm}
Begin the session by welcoming students and briefly explaining that this opening activity is meant to help everyone feel comfortable and to connect course material to what they find important.

Invite each student to share:
\begin{itemize}
    \item Their name
    \item Their favorite number, how many significant figures it is and why its importent to them or where they/or someone can use it
\end{itemize}

If a student is unsure or has not decided on a favorite number, guide them with follow-up prompts such as:
\begin{itemize}
    \item When you think about your daily routine, what number do you see the most (like a room number, bus route, or locker)?
    \item What number do you use most often in your phone or passwords (even if you don’t think about it)?
    \item What number do you check most often each day (time, date, or steps)?
    \item If you had to pick a number for something you like (sports jersey, pizza slices, or a favorite year), which one would you choose?
    \item What number would you use to describe your day right now (like a 1–10 rating)?
\end{itemize}

Start by acknowledging that numbers are already part of everyday life, and that nobody has to be “good at math” to begin learning how scientists use numbers.
Invite students to notice where they already see numbers in the world around them, and emphasize that today’s work is about making those numbers more useful and meaningful, not about memorizing rules.

Encourage them to think of the next activity as a way of practicing a skill—like learning a new language—rather than proving anything about themselves. Then move directly into the task, letting the activity show them how numbers can feel familiar when they are connected to real situations. \\

\hline
\end{tabular}
\newpage












\begin{tabular}{|p{0.9\linewidth}|}
\hline
\begin{tikzpicture}[remember picture,overlay]
\node at (current page.center) {\begin{tikzpicture}[scale=1]
\node[draw, very thick, minimum width=6in, minimum height=9in] (frame) {};
\node[scale=20, anchor=center] (N) at (-2.5,0) {\textbf{N}};
\node[scale=20, anchor=center] (slash) at (0,0) {\textbf{/}};
\node[scale=20, anchor=center] (A) at (2.5,0) {\textbf{A}};
\draw[thick, dashed] (N.east) -- (slash.west);
\draw[thick, dashed] (slash.east) -- (A.west);
\node[rotate=90, scale=1.5] at (-5,0) {Not};
\node[rotate=90, scale=1.5] at (5,0) {Applicable};

\draw[dotted, thick] (-4,-3) -- (-1,-1);
\draw[dotted, thick] (1,-1) -- (4,-3);
\node[scale=1.2] at (-4.2,-3.3) {Unavailable};
\node[scale=1.2] at (4.2,-3.3) {Undefined};
\end{tikzpicture}};
\end{tikzpicture} \\
\hline
\end{tabular}





\newpage
\noindent
{\large \textbf{Checking for Understanding (questions \& answers)}} \\
\begin{tabular}{|p{0.9\linewidth}|}
\hline
\begin{enumerate}
    \item How do you see yourself as someone who uses numbers in everyday life?
    \begin{itemize}
        \item What everyday situations already require you to use numbers or make estimates?
        \item How do you feel when you have to work with numbers in real life?
    \end{itemize}
    \item What strengths do you bring to working with numbers and problem solving?
    \begin{itemize}
        \item Are there tasks you do well that involve logic, patterns, or practical thinking?
        \item What skills do others notice in you that could support quantitative work?
    \end{itemize}
    \item What challenges do you face when math feels unfamiliar or stressful?
    \begin{itemize}
        \item What situations make numbers feel intimidating?
        \item What strategies help you feel more confident when you try something new?
    \end{itemize}
    \item Imagine yourself in a career five to ten years from now.
    \begin{itemize}
        \item What kinds of problems are you solving on a daily basis?
        \item What role do numbers, measurements, or calculations play in that work?
    \end{itemize}
    \item How does learning measurement and math in this course connect to that future?
    \begin{itemize}
        \item Which skills you build now will be useful in many careers?
        \item What makes math feel more approachable when it is tied to real tasks?
    \end{itemize}
    \item If you are unsure about your future path:
    \begin{itemize}
        \item What types of real-world problems do you enjoy figuring out?
        \item What kinds of jobs do you imagine yourself doing even if you don’t know the title yet?
        \item What skills would you like to be confident using in your future?
    \end{itemize}
    \item What does it mean to you to use math like a professional?
    \begin{itemize}
        \item How do professionals approach new problems, ask questions, and learn from mistakes?
        \item What behaviors can you practice now that will help you in any career?
    \end{itemize}
\end{enumerate} \\
\hline
\end{tabular}










\newpage
\noindent
\section*{\underline{Activity 1}}
\begin{tabular}{|p{4.2cm}|p{9.8cm}|}
\hline
\textbf{Collaborative Learning Techniques} & Clusters \\
\hline
\textbf{Learning Strategy} & Vague to Specific, Popcorn \\
\hline
\textbf{Time} & 12:10---12:35 \\
\hline
\textbf{Materials} & Computers, usb-c to hdmi \\
\hline
\textbf{Lecture/Textbook/Old Plans/Slide References} & Chapter 2: Measurements \\
\hline
\end{tabular}

\vspace{0.5cm}

\noindent
{\large \textbf{Instructions (please provide a \hl{DETAILED} activity breakdown below (and/or attached}} \\


\begin{tabular}{|p{0.9\linewidth}|}
\hline
    Have papers with a few of the following vague statements on which are used as popcorn to have students explain their perspective of how measurments might be used in these vague scenarios.

\begin{itemize}
    \item The hallway is very long.
    \item The classroom is kind of cold.
    \item The cup is almost full.
    \item The metal block is heavy.
    \item The mercury in the beaker is very little, but heavy.
    \item The bag of rice is large.
    \item The book is thin.
    \item The bottle is half empty.
    \item The room has too many students, not enough chairs.
    \item The car drove a long distance.
    \item The paper is small.
    \item The pencil is short.
\end{itemize}

\textbf{Instructions:}
\begin{enumerate}
    \item Form a group of 3--4 students.
    \item Your group will receive several everyday statements that describe something in an imprecise way.
    \item For each statement, rewrite it as a chemistry-style measurement by answering the following:
    \begin{itemize}
        \item What quantity is being measured?
        \item What unit should be used?
        \item Is the value exact or measured (inexact)?
        \item What uncertainty would be reasonable?
        \item How many significant figures make sense?
    \end{itemize}
    \item After rewriting the statement, create one related calculation. This may include:
    \begin{itemize}
        \item a unit conversion,
        \item a density calculation,
        \item a temperature conversion, or
        \item writing a value in scientific notation.
    \end{itemize}
    \item Show all work, including units, and apply the correct significant figure rules.
    \item Be prepared to explain how your group chose the unit, level of precision, and number of significant figures.
\end{enumerate}

\textbf{Reminder:}
There may be more than one reasonable answer. Focus on clear reasoning and correct use of measurement language. \\
\hline
\end{tabular}
\newpage

\begin{tabular}{|p{0.9\linewidth}|}
\hline
\begin{tikzpicture}[remember picture,overlay]
\node at (current page.center) {\begin{tikzpicture}[scale=1]
\node[draw, very thick, minimum width=6in, minimum height=9in] (frame) {};
\node[scale=20, anchor=center] (N) at (-2.5,0) {\textbf{N}};
\node[scale=20, anchor=center] (slash) at (0,0) {\textbf{/}};
\node[scale=20, anchor=center] (A) at (2.5,0) {\textbf{A}};
\draw[thick, dashed] (N.east) -- (slash.west);
\draw[thick, dashed] (slash.east) -- (A.west);
\node[rotate=90, scale=1.5] at (-5,0) {Not};
\node[rotate=90, scale=1.5] at (5,0) {Applicable};

\draw[dotted, thick] (-4,-3) -- (-1,-1);
\draw[dotted, thick] (1,-1) -- (4,-3);
\node[scale=1.2] at (-4.2,-3.3) {Unavailable};
\node[scale=1.2] at (4.2,-3.3) {Undefined};
\end{tikzpicture}};
\end{tikzpicture} \\
\hline
\end{tabular}

\newpage

\noindent
{\large \textbf{Checking for Understanding (questions \& answers)}} \\[0.2em]

\begin{tabular}{|p{0.9\linewidth}|}
\hline
\begin{itemize}
    \item What two parts must every measurement include, and why?
    \item How do you decide which unit to use for a measurement?
    \item How do you tell the difference between an exact number and an inexact number?
    \item What does uncertainty mean in a measurement, and how do you record it?
    \item How do you determine how many significant figures are in a number?
    \item How do you decide how many significant figures to keep after multiplying or dividing?
    \item How do you decide how many significant figures to keep after adding or subtracting?
    \item How do you convert between units using conversion factors?
    \item How do you use dimensional analysis to set up a calculation?
    \item How can density be used as a conversion factor?
    \item How do you convert between Celsius, Fahrenheit, and Kelvin?
    \item How do you write a number in scientific notation and convert it back to decimal form?
\end{itemize}
\\
\hline
\end{tabular}
\newpage


\newpage
\noindent
\section*{\underline{Activity 2}}
\begin{tabular}{|p{4.2cm}|p{9.8cm}|}
\hline
\textbf{Collaborative Learning Techniques} & Fishbowl \\
\hline
\textbf{Learning Strategy} & Reverse Engingeering \\
\hline
\textbf{Time} & 12:25---12:40 \\
\hline
\textbf{Materials} & Notes \\
\hline
\textbf{Lecture/Textbook/Old Plans/Slide References} & Chapter 2: Measurements \\
\hline
\end{tabular}

\vspace{0.5cm}
\noindent

\noindent
{\large \textbf{Instructions (please provide a \hl{DETAILED} activity breakdown below (and/or attached)}} \\[0.2em]

\begin{tabular}{|p{0.9\linewidth}|}
\hline
\textbf{Objective:}
Students will develop, test, and defend rules for determining significant figures by engaging in structured peer critique and evidence-based rebuttal.

\subsubsection*{Group Roles}

\begin{itemize}
    \item \textbf{Inner Group (Rule Constructors):} Infers rules for significant figures from numerical examples with known precision and presents a complete rule set.
    \item \textbf{Outer Group (Rule Testers and Challengers):} Independently develops its own rule definitions, tests the inner group’s rules, and defends alternative rules when failures occur.
\end{itemize}

\subsubsection*{Phase 0: Independent Rule Construction (Outer Group)}

Before any numerical examples are discussed, the outer group:
\begin{itemize}
    \item Writes its own definitions for significant figure rules.
    \item Does not apply these rules to the examples yet.
\end{itemize}

\subsubsection*{Phase 1: Evidence-Based Rule Construction (Inner Group)}

The inner group is shown numerical examples where the correct number of significant figures is provided.

For each number, the inner group:
\begin{itemize}
    \item States the given number of significant figures.
    \item Identifies features of the number that justify that count.
    \item Uses these observations to construct a rule set.
\end{itemize}

The outer group listens silently during this phase.\\
\hline
\end{tabular}

\newpage
\noindent
{\large \textbf{Instructions (please provide a \hl{DETAILED} activity breakdown below (and/or attached)}} \\[0.2em]

\begin{tabular}{|p{0.9\linewidth}|}
\hline
\subsubsection*{Phase 2: Rule Presentation (Inner Group)}

The inner group formally presents its complete rule set to the outer group.

No discussion or correction occurs during the presentation.

\subsubsection*{Phase 3: Rule Testing (Outer Group)}

Each number is assigned to one member of the outer group.

For their assigned number, each outer-group member:
\begin{itemize}
    \item Applies the \emph{inner group’s rules} step by step.
    \item States the predicted number of significant figures.
    \item Compares the prediction to the given correct value.
\end{itemize}

\subsubsection*{Phase 4: Rebuttal and Rule Defense}

If the inner group’s rules fail for a number, the outer-group member responsible for that number:
\begin{itemize}
    \item Applies the \emph{outer group’s own rules} to the same number.
    \item Demonstrates why those rules correctly produce the given significant figure count.
    \item Explains which assumption or omission caused the inner group’s rule to fail.
\end{itemize}

The inner group may respond by revising or defending their rule.

\subsubsection*{Whole-Class Synthesis}

As a class:
\begin{itemize}
    \item Compare competing rule sets.
    \item Identify which rules were robust across all cases.
    \item Finalize a complete and correct set of significant figure rules grounded in the evidence.
\end{itemize} \\

\hline
\end{tabular}


\newpage
\noindent
{\large \textbf{Content (fully worked out problems \& answers)}} \\[0.2em]

\begin{tabular}{|p{0.9\linewidth}|}
\hline
\paragraph{Set A: Nonzero Digits and Interior Zeros}

\begin{center}
\begin{tabular}{c c}
\textbf{Number} & \textbf{Significant Figures} \\
\hline
$347$ & $3$ \\
$4.52$ & $3$ \\
$1.007$ & $4$ \\
$10.05$ & $4$
\end{tabular}
\end{center}

\paragraph{Set B: Leading Zeros}

\begin{center}
\begin{tabular}{c c}
\textbf{Number} & \textbf{Significant Figures} \\
\hline
$0.0064$ & $2$ \\
$0.040$ & $2$ \\
$0.000901$ & $3$
\end{tabular}
\end{center}

\paragraph{Set C: Trailing Zeros Without a Decimal Point}

\begin{center}
\begin{tabular}{c c}
\textbf{Number} & \textbf{Significant Figures} \\
\hline
$1500$ & $2$ \\
$12000$ & $2$ \\
$700$ & $1$
\end{tabular}
\end{center}

\paragraph{Set D: Trailing Zeros With a Decimal Point}

\begin{center}
\begin{tabular}{c c}
\textbf{Number} & \textbf{Significant Figures} \\
\hline
$15.00$ & $4$ \\
$120.0$ & $4$ \\
$7.0$ & $2$
\end{tabular}
\end{center}

\paragraph{Set E: Scientific Notation}

\begin{center}
\begin{tabular}{c c}
\textbf{Number} & \textbf{Significant Figures} \\
\hline
$3.20 \times 10^{4}$ & $3$ \\
$6.0 \times 10^{-3}$ & $2$ \\
$1.005 \times 10^{2}$ & $4$
\end{tabular}
\end{center}\\
\hline
\end{tabular}
\newpage

\vspace{0.5cm}
\noindent
{\large \textbf{Checking for Understanding (questions \& answers)}}

\begin{center}
\begin{tabular}{|p{0.47\linewidth}|p{0.47\linewidth}|}
\hline
\textbf{Question} & \textbf{Expected Answer} \\
\hline
What was the goal of the inner group during the activity? &
To infer significant figure rules by analyzing numerical examples where the correct number of significant figures was already given. \\
\hline
What was the goal of the outer group before the inner group presented their rules? &
To independently write their own definitions of significant figure rules based on prior knowledge, without applying them to the examples. \\
\hline
Why was the outer group not allowed to apply any rules during the inner group’s discussion? &
So they could listen objectively, avoid influencing the rule construction, and preserve their own independent rule set for later comparison. \\
\hline
What did the inner group do once they finished reverse-engineering the rules? &
They formally presented their complete rule set to the outer group without interruption. \\
\hline
What was each outer-group member responsible for during rule testing? &
Applying the inner group’s rules to one assigned number and checking whether the predicted significant figure count matched the given value. \\
\hline
What happened if the inner group’s rules did not work for a number? &
The outer-group member responsible for that number applied the outer group’s own rules and explained why those rules produced the correct number of significant figures. \\
\hline
Why was it important that the correct significant figure counts were already provided? &
So the focus remained on identifying patterns and testing rules rather than debating numerical correctness. \\
\hline
What role did disagreement play in the activity? &
Disagreement was used as evidence to refine, revise, or replace incomplete or incorrect rules. \\
\hline
What skill beyond significant figures was this activity designed to practice? &
Constructing claims from evidence, testing models, and defending or revising ideas through peer critique. \\
\hline
What indicates that a significant figure rule is robust? &
It correctly explains all given examples without contradictions or special-case exceptions. \\
\hline
\end{tabular}
\end{center}
\newpage














\newpage
\noindent
\section*{\underline{Closer}}
\begin{tabular}{|p{4.2cm}|p{9.8cm}|}
\hline
\textbf{Collaborative Learning Techniques} & Walkthrough \\
\hline
\textbf{Learning Strategy} & Measure, Listen, Remeasure \\
\hline
\textbf{Time} & 12:40---12:50 \\
\hline
\textbf{Materials} & A computer, textbook, ready to speak \\
\hline
\textbf{Lecture/Textbook/Old Plans/Slide References} &  Chapter 2: Measurements \\
\hline
\end{tabular}

\vspace{0.5cm}
\noindent
{\large \textbf{Instructions (please provide a \hl{DETAILED} activity breakdown below (and/or attached}} \\[0.2em]

\begin{tabular}{|p{0.9\linewidth}|}
\hline
Explain to students that the next activity will focus on identifying how different chemistry ideas feel in terms of difficulty. Emphasize that this is about awareness, not correctness.

Introduce each topic briefly by describing the \emph{type of idea} it addresses without giving definitions, facts, terminology, or examples.

\paragraph{Topic Introductions}
As you move through the content, describe each topic using broad language such as:

\begin{itemize}
    \item How do we describe a measurement, and why must every measurement include both a number and a unit?
    \item How do the English and metric systems compare, and which system do we use most often in chemistry?
    \item How do we choose and use the correct metric base units for length, mass, and volume?
    \item How do metric prefixes change the size of a unit, and how do they help us work with very large or very small numbers?
    \item How do we tell the difference between exact numbers and measured (inexact) numbers?
    \item How do we record a measurement correctly, including the certain digits and the first uncertain digit?
    \item How do we determine the number of significant figures in a given value?
    \item How do we apply significant figure rules when adding, subtracting, multiplying, or dividing values?
    \item How do we use conversion factors, dimensional analysis, density, and temperature scales to solve chemistry measurement problems?
\end{itemize}


Avoid answering any questions at this stage. \\
\hline
\end{tabular}
\newpage

\vspace{0.5cm}
\noindent
{\large \textbf{Instructions (please provide a \hl{DETAILED} activity breakdown below (and/or attached}} \\[0.2em]

\begin{tabular}{|p{0.9\linewidth}|}
\hline

Create the following headings on the board. Under each heading, list each student’s difficulty score (1--10).

\begin{enumerate}
    \item Measurements
    \item Systems of Measurements
    \item Metric system units
    \item Common metric system prefixes
    \item Exact numbers
    \item Inexcat numbers
    \item Uncertainty in measurements
    \item Significant figures
    \item Significant figures and mathematical operations
    \item Scientific notation
    \item Conversion factors
    \item Dimensional analysis
    \item Density
    \item Temperature
\end{enumerate} \\
\hline
\end{tabular}
\newpage
{\large \textbf{Instructions (please provide a \hl{DETAILED} activity breakdown below (and/or attached}} \\[0.2em]

\begin{tabular}{|p{0.9\linewidth}|}
\hline
After introducing each heading, ask students to provide a number from 1 to 10 indicating how challenging the concept feels to them. Record all numbers before moving on.

Once all difficulty scores are recorded, explain that these numbers form a dataset.

Walk through the averaging process slowly and deliberately on the board.

\paragraph{Averaging Walkthrough}
For a single topic:

\begin{enumerate}
    \item Count how many scores were recorded.
    \item Add all the values to get a total.
    \item Choose a reasonable whole-number estimate for the average.
    \item Multiply that estimate by the number of measurements.
    \item Compare the result to the total sum.
\end{enumerate}

Explain aloud whether the estimate is too high or too low and why.

Adjust the estimate and repeat until the value is close to the total sum.

Round to the nearest whole number and label it clearly as the average difficulty for that topic.

\begin{itemize}
    \item These numbers function like repeated measurements of the same quantity.
    \item Averaging helps represent typical behavior.
    \item This same reasoning is used when averaging measurements from physical samples.
\end{itemize} \\
\hline
\end{tabular}
\newpage








\noindent
{\large \textbf{Content (fully worked out problems \& answers)}} \\[0.2em]

\begin{tabular}{|p{0.9\linewidth}|}
\hline
\textbf{(1)}. Define a measurement. Provide examples of measurable phenomena.\\
\textbf{Solution:} \\
Measurements are determinations of the dimensions, capacity, quantity, or extent of something. \\
A few examples of measurable phenomena: \\
\begin{itemize}
    \item Length
    \item Time
    \item Area
    \item Volume
    \item Speed
    \item Temperature
    \item Density
\end{itemize} \\
\hline
\end{tabular}







\begin{tabular}{|p{0.9\linewidth}|}
\hline
\textbf{(2)} What is the difference between SI (Système International d'Unités) and English in terms of measurements?\\
\textbf{Solution:} SI originated from France and is a simplified set of definitions which relate measurements by a factor of $10^x$. \\
English is a system of measurments which originated from England which relate measurment by a factor other than $10^x$. \\
\hline
\end{tabular}
















\begin{tabular}{|p{0.9\linewidth}|}
\hline
\textbf{(3)} Define a base unit. Provide base units examples.\\
\textbf{Solution:} Specifically in SI, before any further measurments are taken in reference to another, a standard is measured and called the base unit. \\
The base unit is primitive, stored and authoritative \- its not derived from anything within the system. \\
\begin{itemize}
    \item gram
    \item inch
    \item liter
    \item meter
    \item kelvin
\end{itemize} \\
\hline
\end{tabular}
\newpage

\noindent
{\large \textbf{Content (fully worked out problems \& answers)}} \\[0.2em]

\begin{tabular}{|p{0.9\linewidth}|}
\hline
\textbf{(4)} List common english units. \\
\textbf{Solution:} \\
\begin{itemize}
    \item inch
    \item foot
    \item quart
    \item gallon
    \item mile
    \item miles per hour
    \item pound
    \item stone
    \item yard
    \item cup
    \item ounce
    \item fahrenheit
\end{itemize} \\
\hline
\end{tabular}



\begin{tabular}{|p{0.9\linewidth}|}
\hline
\textbf{(5)} Define an SI derived unit. Explain what a prefix means in this context. Provide definitions of some dervied units. \\
\textbf{Solution:} SI derived units are relationships between a base unit and another unit, or between two prefixed units which are commonly represented by an equality definition \\
\begin{itemize}
    \item $1~g$ (gram) $=~1000~mg$ (milligrams)
    \item $10~mg$ (milligrams) $=~1~cg$ (centigram)
    \item $1000~m$ (meters) $=~1~km$ (kilometer)
    \item $1,000,000~g$ (grams) $=~1~Mg$ (megagram)
    \item $10 \frac{m}{s}$ (meters per second) $=~1~\frac{dam}{s}$ (decameters per second)
    \item $1000~mL$ (milliliters) $= 1~L$ (Liter)
    \item $1~cm^3$ (cubic centimeters) $= 1~mL$ (milliliter) $=~1~cc$ (cubic centimeters)
    \item $ 273.15~K$ (kelvin) $=~0~\circ\mathrm{C}$ (centigrade)
\end{itemize} \\
\hline
\end{tabular}




\newpage
\noindent
{\large \textbf{Content (fully worked out problems \& answers)}} \\[0.2em]

\begin{tabular}{|p{0.9\linewidth}|}
\hline
\textbf{(6)} Define a conversion factor. Provide examples of conversion factors between English and SI units and which measurable phenomena they represent. \\
\textbf{Solution:} \\
\begin{itemize}
    \item $32~^\circ\mathrm{F}$ (fahrenheit) $=~0~^\circ\mathrm{C}$ (centigrade)
    \item $1~in$ (inch) $=~2.54~cm$ (centimeters) | Length
    \item $1~ft$ (foot) $=~0.3048~m$ (meters) | Length
    \item $1~in^2$ (square inch) $=~6.4516~cm^2$ (square centimeters) | Area
    \item $1~acre$ $=~4046.86~m^2$ (square meters) | Area
    \item $1~gallon$ $=~3.78541~L$ (liters) | Volume
    \item $1~psi$ (pounds per square inch) $=~6894.76~Pa$ (pascals) | Pressure
    \item $1~atm$ (atmosphere) $=~101,325~Pa$ (pascals) | Pressure
    \item $1~pounds$ $=~0.453592~kg$ (killigram) | Mass
    \item $1~ounce$ $=~28.3495~g$ (gram) | Mass
\end{itemize} \\
\hline
\end{tabular}
\newpage


\noindent
{\large \textbf{Content (fully worked out problems \& answers)}} \\[0.2em]

\textbf{(7)} Define an exact number. Explain why counting results in an an exact number. \\
\textbf{Solution:} An exact number is a counted quantity, or a measurment which is used as a standard. Exact numbers have infinite significant figures. \\
\begin{table}
    \centering
    \begin{minipage}[t]{0.48\linewidth}
        \centering

        % Length Conversion Table
        \begin{tabular}{|c|}
            \hline
            \textbf{Length Conversion Factors} \\
            \hline
            $1 \text{ Angstrom ($\AA$)} = 0.1 \text{ nanometers (nm)}$ \\
            $1 \text{ inch (in)} = 2.54 \text{ centimeters (cm)}$ \\
            $1 \text{ foot (ft)} = 0.3048 \text{ meters (m)}$ \\
            $1 \text{ mile (mi)} = 1.60934 \text{ kilometers (km)}$ \\
            \hline
        \end{tabular}
        \caption{Length Conversion Factors}
        \label{tab:length_conversions}
        \vspace{1cm}
        % Volume Conversion Table
        \begin{tabular}{|c|}
            \hline
            \textbf{Volume Conversion Factors} \\
            \hline
            $1 \text{ teaspoon (tsp)} = 4.92892 \text{ milliliters (mL)}$ \\
            $1 \text{ tablespoon (tbsp)} = 14.7868 \text{ milliliters (mL)}$ \\
            $1 \text{ fluid ounce (fl oz)} = 29.5735 \text{ milliliters (mL)}$ \\
            $1 \text{ cup} = 0.236588 \text{ liters (L)}$ \\
            $1 \text{ pint (pt)} = 0.473176 \text{ liters (L)}$ \\
            $1 \text{ quart (qt)} = 0.946353 \text{ liters (L)}$ \\
            $1 \text{ gallon (gal)} = 3.78541 \text{ liters (L)}$ \\
            \hline
        \end{tabular}
        \caption{Volume Conversion Factors}
        \label{tab:volume_conversions}
        \vspace{1cm}
        % Pressure Conversion Table
        \begin{tabular}{|c|}
            \hline
            \textbf{Pressure Conversion Factors} \\
            \hline
            $1 \text{ atmosphere (atm)} = 101.325 \text{ kilopascals (kPa)}$ \\
            $1 \text{ atmosphere (atm)} = 760 \text{ millimeters of Hg (mmHg)}$ \\
            \hline
        \end{tabular}
        \caption{Pressure Conversion Factors}
        \label{tab:pressure_conversions}
        \vspace{1cm}
        % Energy Conversion Table
        \begin{tabular}{|c|}
            \hline
            \textbf{Energy Conversion Factors} \\
            \hline
            $1 \text{ calorie (cal)} = 4.184 \text{ joules (J)}$ \\
            $1 \text{ kilocalorie (kcal)} = 4184 \text{ joules (J)}$     \\
            \hline
        \end{tabular}
        \caption{Energy Conversion Factors}
        \label{tab:energy_conversions}
    \end{minipage}\hfill
    \begin{minipage}[t]{0.48\linewidth}
        % Mass Conversion Table
        \centering
        \begin{tabular}{|c|}
            \hline
            \textbf{Mass Conversion Factors} \\
            \hline
            $1 \text{ unified atomic mass unit (amu)} = 1 \text{ dalton (Da)}$ \\
            $1 \text{ Dalton (Da)} = 1.660054 \times 10^{-21} \text{ milligrams (mg)}$ \\
            $1 \text{ ounce (oz)} = 28.3495 \text{ grams (g)}$ \\
            $1 \text{ pound (lb)} = 0.453592 \text{ kilograms (kg)}$ \\
            $1 \text{ tonne (ton)} = 2000 \text{ pounds (lb)}$ \\
            \hline
        \end{tabular}
        \caption{Mass Conversion Factors}
        \label{tab:mass_conversions}
        \vspace{1cm}
        % Metric Prefixes Table
        \begin{tabular}{|c|c|c|}
            \hline
            \textbf{Prefix} & \textbf{Symbol} & \textbf{Factor} \\
            \hline
            Tera & T & $10^{12}$ \\
            Giga & G & $10^9$ \\
            Mega & M & $10^6$ \\
            Kilo & k & $10^3$ \\
            Hecto & h & $10^2$ \\
            Deka & da & $10^1$ \\
            \hline
            (Base Unit) &  & $10^0$ \\
            \hline
            Deci & d & $10^{-1}$ \\
            Centi & c & $10^{-2}$ \\
            Milli & m & $10^{-3}$ \\
            Micro & $\mu$ & $10^{-6}$ \\
            Nano & n & $10^{-9}$ \\
            Pico & p & $10^{-12}$ \\
            Femto & f & $10^{-15}$ \\
            Atto & a & $10^{-18}$ \\
            \hline
        \end{tabular}
        \caption{Common Metric Prefixes}
        \label{tab:metric_prefixes}
        \vspace{1cm}
        % Time Conversion Table
        \begin{tabular}{|c|c|}
            \hline
            \textbf{Time Unit Conversion Factors} \\
            \hline
            $1 \text{ minute (min)} = 60 \text{ seconds (s)}$ \\
            $1 \text{ hour (hr)} = 60 \text{ minutes (min)}$ \\
            $1 \text{ day (day)} = 24 \text{ hours (hr)}$ \\
            $1 \text{ week (wk)} = 7 \text{ days (day)}$ \\
            $1 \text{ month (mo)} \approx 30.44 \text{ days (day)}$ \\
            $1 \text{ year (yr)} \approx 365.25 \text{ days (day)}$ \\
            \hline
        \end{tabular}
        \caption{Time Conversions}
        \label{tab:time_conversions}
    \end{minipage}
\end{table}

\newpage
\noindent
{\large \textbf{Content (fully worked out problems \& answers)}} \\[0.2em]

\begin{tabular}{|p{0.9\linewidth}|}
\hline
\textbf{(8)} Define an inexact number. Explain the difference between precision and accuracy. Explain the important of calibration and user repeatability. \\
\textbf{Solution:} \\
An inexact number is a measurement with a degree of uncertainty. \\
Every measurement device varies in its degree of precision, which results in an inexact number with varying uncertainty. \\
A measurment device must be calibrated to match as close to the exact value as possible. \\
\hline
\end{tabular}




\begin{tabular}{|p{0.9\linewidth}|}
\hline
\textbf{(9)} Define significant figures. What do they mean? How are they determined from an exact number? An inexact number? \\
\textbf{Solution:} Significant figures are determined by a set of rules about an exact or inexact number. They are representative of precision. \\
An exact number has infinite significant figures.
An inexact number follows a few rules: \\
\begin{enumerate}
    \item Leading zeroes are insignificant, regardless of a decimal point.
    \item Numeric digits which are not zero are significant.
    \item Zeroes between numeric digits or a decimal point are significant.
    \item Trailing zeroes are always significant after the decimal point.
\end{enumerate} \\
\hline
\end{tabular}




\begin{tabular}{|p{0.9\linewidth}|}
\hline
\textbf{(10)} Determine the significant figures of these values: \\
\textbf{Solution:} \\
\begin{itemize}
    \item $100$ = 1 significant figure
    \item $100.$ = 3 significant figures
    \item $12$ donuts in a dozen = infinite significant figures
    \item $808~students$ = infinite significant figures
    \item $808~miles$ = 3 significant figures
    \item $0.000432~ft$ = 3 significant figures
    \item $1.002~in$ = 4 significant figures
    \item $2000~lbs$ in a ton = infinite significant figures
    \item $3.350~g$ = 4 significant figures
\end{itemize} \\
\hline
\end{tabular}





\newpage
\noindent
{\large \textbf{Content (fully worked out problems \& answers)}} \\[0.2em]

\begin{tabular}{|p{0.9\linewidth}|}
\hline
\textbf{(11)} Determine the minimum and maximum significant figures of these measurement devices. \\
\textbf{Solution:} \\
\begin{itemize}
    \item A $1~ft$ ruler with increments of $0.1~in$ = maximum of 5 significant figures ($11.015~in$), minimum of 1 significant figure ($0.005~in$).
    \item A digital scale which measures up to $500~g$ with increments to the nearest $0.001~g$ = maximum of 6 significant figures ($400.006~g$), a minimum of 1 significant figure ($0.006~g$).
    \item An analog speedometer which measures up to $120~mph$ with increments of $1~mph$ = maximum of 4 significiant figures ($100.5~mph$), minimum of 1 significant figure ($0.5~mph$).
    \item A digital speedometer which measures up to $105~mph$ with increments to the nearest $0.1~mph$ = maximum of 4 significant figures ($102.1~mph$), minimum of 1 significant figure ($0.1~mph$).
    \item A graduated cylinder which measures up to $100~mL$ with increments to the nearest $1~mL$ = maximum of 3 significant figures ($89.5~mL$), minimum of 1 significant figure ($0.5~mL$).
    \item A gallon ($1~gal$) with no increments = maximum of 2 significant figures $0.3~gal$, minimum of 1 significant figure ($0.4~gal$).
    \item An $12~hr$ analog clock with increments to the nearest $1~min$ = maximum of 4 significant figures for hours ($12.02~hr$), 3 significant figures for minutes $35.5~min$, 1 significant figure for seconds $30 sec$), minimum of 1 significant figure ($10.0~hr$) | ($0.5~min$) | ($30~sec$).
    \item A $24~hr$ digital clock with increments to the nearest $1~sec$ = maximum of 6 significant figures ($23.0202~hr$) and minimum of 1 significant figure ($1~sec$).
\end{itemize} \\
\hline
\end{tabular}





\begin{tabular}{|p{0.9\linewidth}|}
\hline
\textbf{(12)} What is scientific notation and how is it used? \\
\textbf{Solution:} Scientific notation uses base 10 to multiply a value by to reach the actual value. \\
Scientific notation is used to condense rather large numbers with many digits but not many significant figures into a more palateable value. \\
Scientific notation allows for rounding to the nearest 1/10th of a digit and can be used to represent an inexact number with the correct significant figures. \\
\begin{itemize}
  \item What is $0.004560$ in scientific notation? --- $4.560 \times 10^{-3}$
  \item Express $72{,}300$ with correct significant figures. --- $7.23 \times 10^{4}$
  \item Convert $6{,}022{,}000{,}000{,}000{,}000{,}000{,}000$ to scientific notation. --- $6.022 \times 10^{23}$
  \item What is $9.80 \times 10^{-2}$ in decimal form? --- $0.0980$
  \item Write $1.000 \times 10^{3}$ in standard notation. --- $1000$
  \item Express $0.000120$ in scientific notation. --- $1.20 \times 10^{-4}$
\end{itemize}
\hline
\end{tabular}
\newpage




\noindent
{\large \textbf{Content (fully worked out problems \& answers)}} \\[0.2em]

\begin{tabular}{|p{0.9\linewidth}|}
\hline
\textbf{(13)} How are significant figures determined of a value returned from a calculation? \\
\textbf{Solution:} For addition and subtraction, the value with the lowest number of significant figures past the decimal point is the limiting value. \\
The caculated value will be rounded to have the same number of significant figures past the decimal point. \\
For multiplication and division, the value with the lowest total numbe rof significant figures in the limiting value. \\
The calculated value will be rounded to have the same number of significant figures as the limiting value.
\begin{itemize}
  \item What is $2.50 \times 3.0$ with correct significant figures? Answer: $2.50 \times 3.0 = 7.5$ (two significant figures).
  \item If a rod is measured as $12\ \mathrm{cm}$ exactly cut in half, what is each half? Answer: $\frac{12\ \mathrm{cm}}{2} = 6.0\ \mathrm{cm}$ (2 is exact).
  \item What is $4.56 + 1.2$ with proper precision? Answer: $4.56 + 1.2 = 5.8$.
  \item A mass of $2.345\ \mathrm{kg}$ is tripled exactly; what is the result? Answer: $3 \times 2.345 = 7.035\ \mathrm{kg}$.
  \item What is $6.022 \times 10^{23} \times 2.0$ with correct significant figures? Answer: $1.2 \times 10^{24}$.
  \item Divide $9.60$ by $3.00$. Answer: $\frac{9.60}{3.00} = 3.20$.
  \item What is $1.000 \times 10^{-3} + 2.5 \times 10^{-4}$? Answer: $1.25 \times 10^{-3}$.
  \item A count of $24$ samples has an average mass of $0.250\ \mathrm{g}$; what is the total mass? Answer: $24 \times 0.250 = 6.00\ \mathrm{g}$ (24 is exact).
\end{itemize}
\hline
\end{tabular}





\begin{tabular}{|p{0.9\linewidth}|}
\hline
\textbf{(14)} What is dimensional analysis and how is it performed? \\
\hline
\textbf{Solution:} \\
\begin{itemize}
    \item Convert 125 miles into kilometers. (Use $1~\text{mile} = 1.609~\text{km}$)
    \item A car travels 95.0 kilometers in 1.20 hours. What is its speed in meters per second?
    \item The density of aluminum is $2.70~\text{g/cm}^3$. What is the mass of a block of aluminum with a volume of $15.0~\text{cm}^3$?
    \item A patient’s mass is 154 lb. Convert this to kilograms. (Use $1~\text{lb} = 453.6~\text{g}$)
    \item Convert 0.750 L into cubic centimeters. (Use $1~\text{L} = 1000~\text{cm}^3$)
    \item A cube has an edge length of 2.50 cm. What is its volume in cubic meters?
    \item The density of mercury is $13.6~\text{g/mL}$. What is the volume (in mL) of a 125 g sample of mercury?
    \item Convert 3.25 hours to seconds.
    \item Convert 48.0 inches to centimeters. (Use $1~\text{in} = 2.54~\text{cm}$)
    \item Convert 2.50 gallons to liters. (Use $1~\text{gal} = 3.785~\text{L}$)
    \item A metal sample has mass $12.3~\text{g}$ and density $8.90~\text{g/cm}^3$. What is its volume in mL?
    \item Convert $7.50~\text{m/s}$ to $\text{km/h}$.
    \item A sample has mass $27.0~\text{g}$ and volume $35.0~\text{mL}$. Find its density in $\text{g/mL}$.
    \item Convert $745~\text{torr}$ to atm. (Use $1~\text{atm}=760~\text{torr}$
    \item A pump delivers 125.0 L/min. Express the flow rate in $\text{m}^3/\text{s}$.
    \item Convert 512~torr to kPa. (Use $760~\text{torr}=101.325~\text{kPa}$
    \item Convert 1875~g to pounds. (Use $1~\text{lb}=453.6~\text{g}$)
    \item Convert blood glucose level of 125.0 mg/dL to g/L.
    \item Convert 88.0 kPa to atm. (Use $1~\text{atm}=101.325~\text{kPa}$)
    \item A medication dose is 75.0 mg for a 62.0 kg patient. Express this as ppm (mg/kg).
    \item A true density is 8.26 g/mL; a student measures 8.10 g/mL. What is the percent error?
    \item Show that 2.75 L equals 2.75 dm$^3$ using unit factors.
    \item Convert 987654 s to days.
    \item Convert 45.0 grains to grams. (Use $1~\text{grain}=64.79891~\text{mg}$)
    \item Convert 85.0 kg to stones. (Use $1~\text{lb}=0.45359237~\text{kg}$, $1~\text{stone}=14~\text{lb}$)
\end{itemize}
\hline
\end{tabular}


\newpage
\noindent
{\large \textbf{Checking for Understanding (questions \& answers)}} \\[0.2em]

\begin{tabular}{|p{0.9\linewidth}|}
\hline
Remeasure the student's difficulty of each subject and go through averaging again, taking volunteers if they want to do it.
\begin{enumerate}
    \item While we worked through the activity, what did you notice about how you were taking in new information?
    \begin{itemize}
        \item Did you rely more on listening, writing, watching, or discussing?
        \item At what points did something begin to make more sense?
    \end{itemize}

    \item What thinking strategies did you try when you were unsure?
    \begin{itemize}
        \item Did you make estimates, comparisons, or educated guesses?
        \item Did you adjust your thinking as new information appeared?
    \end{itemize}

    \item During the walkthrough, what skills did you actively practice?
    \begin{itemize}
        \item Interpreting information step-by-step
        \item Checking whether a result made sense
        \item Revising an initial idea when it did not fit
    \end{itemize}

    \item How did assigning a number to difficulty change how you thought about the topic?
    \begin{itemize}
        \item Did it make the challenge feel more manageable?
        \item Did it help separate confusion from uncertainty?
    \end{itemize}

    \item Why might it be useful to measure how hard a concept feels before studying it deeply?
    \begin{itemize}
        \item How could this help you estimate how long to spend on a topic?
        \item How could it guide which topics to start with?
    \end{itemize}

    \item When one idea builds on another, how might difficulty ratings help you plan your learning?
    \begin{itemize}
        \item What happens if a foundational idea feels unclear?
        \item How might you adjust your approach when topics are connected?
    \end{itemize}

    \item How does this process compare to how measurements are used in science?
    \begin{itemize}
        \item Why do scientists take multiple measurements?
        \item How does averaging help represent what is typical rather than extreme?
    \end{itemize}

    \item Going forward, how could you apply this way of thinking to studying outside of this session?
    \begin{itemize}
        \item How might you track difficulty over time?
        \item How could this help you study more intentionally?
    \end{itemize}
\end{enumerate} \\
\hline
\end{tabular}

\newpage
\section*{Plan Checklist}

\begin{tabular}{|p{0.9\linewidth}|}
\hline
\begin{itemize}
\item[$\Box$] Variety of Collaborative Learning Techniques and Learning Strategies
\item[$\Box$] Chapter/slide \# where information is coming from
\item[$\Box$] Included timestamps (and buffer time as needed)
\item[$\Box$] Details on how leader will break up students
\item[$\Box$] Checking for Understanding questions
\item[$\Box$] Fully worked out problems and examples for each activity
\item[$\Box$] Necessary links required for online implementation (if applicable)
\item[$\Box$] Source material correctly cited (if applicable)
\item[$\Box$] Plan is uploaded to Teams
\end{itemize}
\\ \hline
\end{tabular}

\end{document}
